%%%%%%%%%%%%%%%%%%%%%%%%%%%%%%%%%
%
%       EXERCÍCIO
%
%%%%%%%%%%%%%%%%%%%%%%%%%%%%%%%%%

\ifdefstring{\atividade}{prova}{%
    \renewcommand{\valorquestao}{\ValorQ\ pontos}
}%

\begin{exercicioBanco}[\valorquestao]
    Seja \(X\) uma variável aleatória com função de probabilidade dada na tabela a seguir:
    \[
        \begin{array}{c|c|c|c|c|c}
            x & 1 & 2 & 3 & 4 & 5 \\
            \hline
            f_{X}(x) & p^{2} & p^{2} & p & p & p^{2}
        \end{array}
    \]
    \begin{enumerate}[(a)]
        \item Encontre o valor de \(p\) para que \(f_{X}(x)\) seja, de fato, uma função de probabilidade.
        %
        \item Calcule \(E(X)\) e \(V(X)\).
    \end{enumerate}
\end{exercicioBanco}

